% Vietnamese translation for OI Public Library
% Source-File: IOI中国国家候选队论文/2005/胡伟栋/附件/polygon-1.2.rtf
\documentclass[11pt]{article}
\usepackage{oipl-vn}

\title{IOI 2004: Polygon\\{\large Phát biểu bài toán}}
\author{IOI 2004, Athens}
\date{Phiên bản 1.1}

\begin{document}
\maketitle

\origtitle{Polygon}
\sourcepath{IOI Candidate Team Papers/2005/Hu Weidong/appendix/polygon-1.2.rtf}

\section{Bài toán}

Một đa giác gồm tất cả các điểm nằm trên biên hoặc nằm bên trong biên của nó.
Một đa giác lồi có tính chất sau: với mọi hai điểm \(P\) và \(Q\) thuộc đa
giác, đoạn thẳng nối \(P\) với \(Q\) cũng nằm trong đa giác.

Trong bài này, mọi đa giác đều là đa giác lồi có ít nhất hai đỉnh, mọi đỉnh
đôi một khác nhau và có tọa độ nguyên. Không có ba đỉnh nào thẳng hàng. Từ
``đa giác'' trong phần còn lại luôn được hiểu theo nghĩa này.

Cho hai đa giác \(A\) và \(B\), \term{tổng Minkowski} của \(A\) và \(B\) gồm
tất cả các điểm có dạng
\[
  (x_1+x_2,\ y_1+y_2),
\]
trong đó \((x_1,y_1)\) là một điểm của \(A\), còn \((x_2,y_2)\) là một điểm
của \(B\). Tổng Minkowski của hai đa giác cũng là một đa giác.

Bài toán xét phép ngược lại. Với một đa giác \(P\), cần tìm hai đa giác
\(A\) và \(B\) sao cho:
\begin{itemize}
  \item \(P\) là tổng Minkowski của \(A\) và \(B\);
  \item \(A\) có từ \(2\) đến \(4\) đỉnh khác nhau, tức \(A\) là một đoạn
  thẳng, một tam giác hoặc một tứ giác;
  \item số đỉnh của \(A\) phải lớn nhất có thể: nếu có thể thì \(A\) là tứ
  giác; nếu không thể là tứ giác thì \(A\) là tam giác nếu có thể; nếu không
  thì \(A\) là đoạn thẳng.
\end{itemize}

Rõ ràng cả \(A\) lẫn \(B\) đều không thể bằng \(P\), vì khi đó đa giác còn
lại sẽ phải là một điểm, mà một điểm không phải là đa giác hợp lệ.

Bạn được cho một tập các tệp vào, mỗi tệp mô tả một đa giác \(P\). Với mỗi
tệp vào, hãy tìm các đa giác \(A\) và \(B\) theo yêu cầu trên, rồi tạo một
tệp ra mô tả \(A\) và \(B\). Với các tệp vào đã cho, luôn tồn tại các đa giác
\(A\) và \(B\) như vậy. Nếu có nhiều kết quả đúng, chỉ cần xuất một kết quả.
Bạn không cần nộp chương trình, chỉ cần nộp các tệp kết quả.

\section{Dữ liệu vào}

Có \(10\) bộ dữ liệu trong các tệp văn bản tên là
\texttt{polygon1.in}, \texttt{polygon2.in}, \(\ldots\),
\texttt{polygon10.in}. Mỗi tệp được tổ chức như sau.

Dòng đầu tiên chứa một số nguyên \(N\): số đỉnh của đa giác \(P\). \(N\) dòng
tiếp theo mô tả các đỉnh theo thứ tự ngược chiều kim đồng hồ, mỗi dòng một
đỉnh. Dòng \(i+1\), với \(i=1,2,\ldots,N\), chứa hai số nguyên \(X_i\) và
\(Y_i\), cách nhau bởi một dấu cách, là tọa độ của đỉnh thứ \(i\). Mọi tọa độ
trong dữ liệu vào đều là số nguyên không âm.

\section{Dữ liệu ra}

Cần nộp \(10\) tệp ra tương ứng với các tệp vào đã cho. Mỗi tệp ra mô tả hai
đa giác \(A\) và \(B\).

Dòng đầu tiên phải chứa đúng văn bản
\begin{verbatim}
#FILE polygon I
\end{verbatim}
trong đó số nguyên \(I\), \(1\le I\le 10\), là số hiệu của tệp vào tương ứng.

Định dạng còn lại tương tự dữ liệu vào. Dòng thứ hai chứa số nguyên
\(N_A\), là số đỉnh của \(A\), với
\[
  2\le N_A\le 4.
\]
\(N_A\) dòng tiếp theo mô tả các đỉnh của \(A\) theo thứ tự ngược chiều kim
đồng hồ. Dòng \(i+2\), với \(i=1,2,\ldots,N_A\), chứa hai số nguyên \(X\) và
\(Y\), cách nhau bởi một dấu cách, là tọa độ của đỉnh thứ \(i\) của \(A\).

Dòng \(N_A+3\) chứa số nguyên \(N_B\), là số đỉnh của \(B\), với
\[
  2\le N_B.
\]
\(N_B\) dòng tiếp theo mô tả các đỉnh của \(B\) theo thứ tự ngược chiều kim
đồng hồ. Dòng \(N_A+j+3\), với \(j=1,2,\ldots,N_B\), chứa hai số nguyên
\(X\) và \(Y\), cách nhau bởi một dấu cách, là tọa độ của đỉnh thứ \(j\) của
\(B\).

\section{Ví dụ}

Tệp ví dụ trong nguồn là \texttt{polygon0.in}:
\begin{verbatim}
5
0 1
0 0
2 0
2 1
1 2
\end{verbatim}

Một tệp ra hợp lệ là:
\begin{verbatim}
#FILE polygon 0
3
0 0
1 0
1 1
3
0 1
0 0
1 0
\end{verbatim}

Với ví dụ này, \(A\) là một tam giác và không thể chọn \(A\) là tứ giác. Tệp
nguồn cũng cho biết có nhiều kết quả đúng; nhiệm vụ chỉ yêu cầu xuất một kết
quả hợp lệ.

\end{document}
